\documentclass[preamble=development-preamble.tex]{markdown}
title: Shared infrastructure for interpreting arithmetic R
date: October 20, 2026

:::Abstract
The library `Interpretation_of_R` provides the proof infrastructure used by
the independent relational and functional interpretations of arithmetic in
the original rank-based pure type system. Its main contributions are a
sound native checker for annotated terms, structural lemmas for dependent
telescopes, fixed-answer negative logic, finite selector observations with
internal domain certificates, and constructive selection of witnesses from
proved inhabitation. This document explains the contracts between these
components and the proof steps that keep finite computation, internal
inhabitation, and source arithmetic distinct. The arithmetic interpretation
theorems themselves belong to the two child libraries.
:::

# Purpose and mathematical boundary

Both arithmetic developments construct terms of the existing calculus.
The syntax is `PTSrk.Syntax.Term nat n`, with finite de Bruijn variables and
annotations on every lambda. The judgment is `PTSrk.Typing.typ` instantiated
with `PTSrk.Ranks.RankSorts`. Writing its sorts as $U_i$, the defining rules
are

$$
 U_i:U_0,
 \qquad
 \frac{\Gamma\vdash A:U_a\quad\Gamma,x:A\vdash B:U_b}
      {\Gamma\vdash\Pi x:A.B:U_{\max(a+1,b)}}.
$$

The library does not introduce arithmetic primitives, a new typing relation,
cumulativity, or eta conversion. `Core.Typed` is an abbreviation for the
original judgment, and `Core.Realizes A` means that a closed original term
has type $A$. The wrappers in `DirectRKernel.v` expose the original typing
constructors and beta-congruence lemmas without changing their premises.

The two clients share these tools but have distinct carriers and equality
relations. The relational library establishes `Theorem4.theorem4`; the
functional library establishes `Theorem7.theorem7` and its effective witness
interfaces. Neither imports the other. Neither depends on either set-model
development. In particular, the common source countermodel discussed below
is not a semantics of the PTS.

# Constructing and checking annotated terms

## Named builders and structural transport

`DirectRCommonCodes.v` provides a small named construction language.
Its constructors retain sort labels, product domains, and lambda
annotations. Assembly resolves names against an explicit environment and
returns an optional scoped original term. `close` additionally requires a
successful closedness test. The named language is therefore a construction
aid; every exported proof concerns the assembled original syntax.

The same module defines the answer type, truth type, arrows, negative
connectives, selectors, and a family of inhabited rank-raising types.
`DirectRStructural.v` proves that weakening, renaming, substitution, and
closed-term lifting interact with these constructions as expected. These
identities matter because a proof assembled under one sequence of binders
must still carry the correct annotations when moved under another.

Dependent observations are represented as telescopes. Each telescope field
records an annotation in the context of all preceding fields. The
formation and abstraction theorems traverse those fields and compute the
product rank. `TelescopeArguments.arguments_application` then types an
application to an actual list of arguments: each argument is checked
against its annotation after substitution of every preceding argument.
This prevents treating a dependent telescope as a list of unrelated types.
Closed arguments can also be lifted uniformly into an arbitrary context.

## The native checker and its exact contract

`DirectREval.v` implements beta evaluation on the original annotated
syntax. Head evaluation exposes a sort or product without reducing an
unused argument first. Full evaluation also visits product domains and
lambda annotations. Every successful elementary step is proved to be an
original beta step, and successful conversion comparison yields original
definitional equality.

`DirectRCheck.v` infers types by structural recursion on the annotated input
term. A numeric fuel parameter bounds only reduction needed to inspect or
compare types. Its results distinguish acceptance, a diagnostic rejection,
and an undecided result caused by exhausted reduction fuel.

:::Theorem[Native checking soundness]
For every scoped context $\Gamma$, term $t$, expected type $A$, and fuel
bound $f$, acceptance of `RkCheck.run_check f` on these inputs implies
$\Gamma\vdash t:A$ in the original calculus. This is the public theorem
`RkCheck.run_check_sound`.
:::

:::Proof
First establish soundness of head inspection and conversion. Induction on
the input term proves `infer_sound`, assuming an original well-formed
context. The application case is the important step: when an inferred type
reduces to $\Pi x:A.B$, original regularity, subject reduction, and product
inversion supply a formation proof for $A$. The checker may then convert
the argument's inferred type to $A$ and apply the original application
rule. A conversion test alone never supplies type formation.

The context checker separately constructs the original, Type-valued
well-formedness evidence by recursion on the context. Finally `check_sound`
checks formation of the expected type before applying the original
conversion rule. Combining these results proves the acceptance theorem.
:::

The tactic `RkCheck.certify_with` uses this theorem and computes a successful
run. Its numeric fuel is a proof-search resource, not a new calculus rule.
No completeness or non-typability theorem is claimed for the checker’s
rejected or undecided outputs. `DirectRSymbolic.v` retains a separate tactic
that assembles original typing and beta-equivalence proofs for symbolic
indices; it likewise adds no typing rule or trusted conversion oracle.

# Negative logic and positive certificates

The fixed answer and truth types are

$$
 \bot:=\Pi Z:U_0.Z:U_1,
 \qquad \top:=U_0\to U_0:U_1.
$$

Truth has its identity inhabitant. The logical constructions use

$$
 \begin{gathered}
 \neg_c A:=A\to\bot,
 \quad A\wedge_c B:=(A\to B\to\bot)\to\bot,
 \\
 \bigvee_c(A_0,\ldots,A_k):=
 (A_0\to\bot)\to\cdots\to(A_k\to\bot)\to\bot.
 \end{gathered}
$$

`DirectRLogic.v` calls $A$ stable when $\neg_c\neg_c A\to A$ is inhabited.
It supplies stability and elimination lemmas for the encoded negative
connectives. In particular, `NegativeLogic.flat_eliminate` eliminates a
finite disjunction into a stable target. It does not permit elimination
into every positive type. The arithmetic translations use this interface
to realize classical deduction without importing a native classical axiom.

Positive evidence must be retained separately. The shared invariant
constructor has the shape

$$
 I(G,C):=\Pi Z:U_8.(G\to C\to Z)\to Z,
$$

where $G:U_6$ is positive observation-domain evidence and $C:U_8$ is a
negative case statement. `positive_invariant_formation` gives this type
rank ten, and `positive_invariant_intro` constructs it from the two pieces.
In a constructor proof one opens the parent's continuation at a fixed
$Z$, eliminates child invariants at that same $Z$, and separately builds
the parent's positive evidence and negative cases. This is the mechanism
that avoids obtaining positive membership by an illicit classical case
split.

Rank raising is also explicit. Set $J_1=\top$ and $J_{k+1}=J_k\to J_k$;
each has a canonical inhabitant. For $A:U_a$, $a\leq k$, and $k\geq2$,
the type $J_{k-1}\to A$ has exactly rank $k$. Abstraction and evaluation at
the canonical inhabitant map between the original and raised types. In
particular, finite-law elimination uses $J_2\to A$ to obtain an exact
rank-three motive. This construction is needed because there is no
cumulativity rule.

# Finite observations represented inside the calculus

## Selectors, tables, and relation tests

For a positive external size $q$, the selector carrier is

$$
 F_q:=\underbrace{U_1\to\cdots\to U_1}_{q\text{ arguments}}\to U_1:U_1,
 \qquad e_i:=\lambda Z_0\ldots Z_{q-1}.Z_i\quad(i<q).
$$

The number of arguments may grow while the carrier rank remains one:
each domain $U_1$ itself has type $U_0$. `DirectRFiniteTables.v` realizes
unary tables, proves selector application, and derives capped successor
and iteration equations. `DirectRBinaryTables.v` adds nested binary
selection. Its `binary_table_action` establishes the equation at every
valid pair of indices, for any table whose entries stay in range.
`DirectRIterations.v` connects explicit repeated application with the
capped finite values.

`DirectRRelationTables.v` uses a matrix of types to encode a finite truth
table. A true entry is $\bot$ and a false entry is $\top$; the interpreted
relation is the negation of the selected entry. Thus a true relation has
the identity $\bot\to\bot$, whereas a proof of a false relation yields
$\bot$ after application to the fixed inhabitant of $\top$.
`negative_true` and `negative_false` formalize these two directions.
The Boolean Horn-law lemmas turn valid finite implications into original
PTS inhabitants, including branches with false premises.

## Internal membership, not an external exhaustiveness assumption

There is no claim that every raw inhabitant of $F_q$ is a selector.
Instead `DirectRPositive.v` defines positive finite membership, specialized
in `DirectRFiniteObservations.v` to

$$
 \delta_q(u):=\Pi Q:F_q\to U_3.
 Q(e_0)\to\cdots\to Q(e_{q-1})\to Q(u):U_4.
$$

Each selector has the corresponding projection certificate. In the
`FiniteObservations` module, `delta_eliminate` specializes it to an
exact rank-three motive and supplies all its selector branches.
`delta_successor` proves membership closure under capped successor;
`BinaryObservations.delta_binary` proves the analogous result for an
arbitrary bounded binary table. To construct output membership, these
proofs open the output predicate and specialize input membership to that
predicate composed with the table operation.

`DirectRFiniteCertificates.v` packages the initial and successor
certificates with their literal annotations for use in observation
telescopes. `DirectRFiniteUp.v` constructs the upward-closed test used in
relational separation. These are actual PTS terms with original typing
proofs, not just native propositions saying a finite algebra has a law.

:::Theorem[Guarded finite-law realization]
Suppose a $k$-argument predicate on $F_q$ has result rank at most three,
and every tuple of valid selector indices has an original proof of that
predicate. Then the predicate holds for all $k$ arguments carrying
$\delta_q$ certificates. The closed, interleaved-guard version is
`GuardedFiniteLaws.guarded_finite_realized`.
:::

:::Proof
`FiniteLawElimination.delta_eliminate_lower` first raises a lower-rank
motive to rank three and lowers the resulting proof by application to the
inhabitant of $J_2$. The theorem `delta_eliminate_arguments` iterates that
construction over all arguments, reducing the result to the supplied
selector-tuple branches. `guarded_intro` introduces each object and its
membership guard before running the eliminator on the remaining body.
This ordering matters: the entire guarded telescope may have rank five,
which is too large for the rank-three membership eliminator.
:::

The finite branch premise quantifies over all tuples of the given length.
Checking a few concrete examples cannot replace it. Conversely, a native
arithmetic proof about table entries is only a side condition until the
table action, negative relation encoding, and membership elimination
lemmas have produced the required internal inhabitant.

# Effective witnesses and the source countermodel

`DirectRDerivationCertificates.v` defines finite trees encoding exactly
the original typing rules. Reduction certificates enumerate original
one-step successors, including reductions in annotations. Conversion
certificates exhibit two paths to a common reduct; completeness uses the
original common-reduct characterization of beta conversion.

:::Theorem[Exact derivation certificates]
`DerivationCertificates.derive_sound` turns successful certificate
evaluation into original typing. Conversely, `derive_complete` gives a
finite certificate for every original typing derivation.
:::

These certificates support a different task from the native checker.
`DirectRWitnessSearch.v` enumerates all finite certificate trees, scans
them for a requested closed type, and proves that an inhabitation proof
implies eventual success. `WitnessSearch.witness` has the contract

$$
 (\exists p.\ \vdash p:A)\ \longrightarrow\quad
 \{p:\operatorname{Term}(0)\mid\vdash p:A\}.
$$

The result uses constructive search over a decidable numeric predicate,
with termination justified by the already proved existence theorem.
It does not extract an arbitrary witness directly from Rocq's `Prop`,
assume a choice axiom, decide unrestricted inhabitation, or provide a
practical running-time bound. This is the common machinery behind the
functional library's effective proof-term interfaces.

`DirectRCountermodel.v` provides the other shared external construction:
$\mathbb N$ with one extra point $\omega$, implemented as `option nat`.
Successor fixes $\omega$; addition and multiplication return zero when an
input is $\omega$. Order is the usual order on naturals, every natural is
below $\omega$, and $\omega$ is below nothing. Ground arithmetic, bounded
downsets, and comparison with a numeral hold, but order reflexivity fails
at $\omega$. Each child library supplies its own source semantics and
deduction soundness before using this structure to establish
nonfaithfulness. The shared structure alone does not prove a source
non-derivability claim.

# Source organization and discussion

The sources form four useful reading groups.

- `DirectRCommonCodes.v`, `DirectRKernel.v`, `DirectREval.v`,
  `DirectRCheck.v`, and `DirectRSymbolic.v` build and certify original terms.
- `DirectRStructural.v` supplies structural transport, and
  `DirectRTelescopeArguments.v` handles dependent application.
  `DirectRLogic.v` and `DirectRPositive.v` distinguish stable negative
  formulas from positive finite membership.
- The `DirectRFinite*` modules, `DirectRBinaryTables.v`,
  `DirectRBinaryObservations.v`, `DirectRRelationTables.v`,
  `DirectRGuardedFiniteLaws.v`, and `DirectRIterations.v` implement the
  uniform finite-observation constructions.
- The certificate and search modules encode original derivations and
  select effective witnesses.
  `DirectRCountermodel.v` supplies the common external source countermodel.

The main design choice is to keep the output contract at original typing
throughout. The checker, symbolic tactics, numerical arithmetic, and
finite tables all assist proof construction, while their soundness
lemmas expose exactly where a PTS inhabitant is produced. Fixed ranks
control eliminations; enlarging a finite table changes the size of terms,
not the rules of the calculus.

The infrastructure itself asserts neither the arithmetic interpretation
theorems nor external uninhabitability of $\bot$. An internal refutation
is a term of type $A\to\bot$; concluding that $A$ has no closed inhabitant
also needs external nontriviality of the answer type. The two statements
must not be conflated. Similarly, successful witness search is effective
in the constructive sense above, without a useful complexity claim.

# Reproducing the verification

The Dune library name is `Interpretation_of_R`, independently of its
physical directory. The pinned environment is OCaml 4.14.2, Dune 3.23.1,
Rocq 9.2.0, Stdlib 9.1.0, and Equations 1.3.2+9.2. Setup and full local CI
are documented in the repository's `docs/DirectR-CI.md`.

From the repository root, in that configured environment:

```sh
dune build theories/Interpretation_of_R/DirectRCheck.vo
dune build theories/Interpretation_of_R/DirectRWitnessSearch.vo
dune build @theories/Interpretation_of_R/relational/audit
dune build @theories/Interpretation_of_R/functional/audit
```

The two audits check public interfaces and their transitive assumptions.
Allocate a fresh absolute Dune build directory when fresh printed
assumptions are required; an incremental success may reuse cached output.
`bash tools/check-local-ci.sh` additionally runs the complete staged
driver, actionlint, and a separate pristine upstream build.

The retained 14 October layout verification includes a fresh complete
build, exact public types, closed assumption reports for the arithmetic
endpoints, and recursive kernel checks. Its command record is
`docs/direct-r-handoff/verification/theories-layout-2026-10-14/COMMANDS.md`.
The recursive loaded-context inventory includes the unchanged upstream
functional-extensionality dependency; this differs from the closed
transitive reports of the arithmetic endpoints. The live handoff records
newer checks and their exact scope. Old comments saying an early draft was
uncompiled describe its authoring history, not the status of the current
checked declarations.
